\section{Zaključak}

Analiza odgovora na pitanja P9 i P10 pokazuje da ispitanici SMR tehnologiju najčešće povezuju sa sigurnosnim značajkama i kraćim vremenom izgradnje, dok se niži troškovi rjeđe ističu kao glavna prednost. Najčešći odgovor na pitanje P9 bile su poboljšane sigurnosne značajke reaktora (33,86\%), nakon čega slijedi kraće vrijeme izgradnje i puštanja u pogon (29,13\%). Time se potvrđuje da percepcija SMR tehnologije u ovom uzorku nije primarno ekonomska, nego je snažno vezana uz sigurnost i izvedivost tehnologije.

Kod pitanja o državnom subvencioniranju razvoja SMR tehnologije većina ispitanika daje potvrdan odgovor. Subvencije smatra opravdanima 59,01\% ispitanika, dok ih 8,01\% ne smatra opravdanima. Istodobno, 32,98\% ispitanika odgovara \textit{Ne znam}, što pokazuje da uz većinsku potporu postoji i znatan prostor za dodatno informiranje javnosti o ekonomskim, sigurnosnim i energetskim posljedicama razvoja SMR tehnologije.

Statistička analiza pokazala je tri značajne povezanosti: između spola i odgovora na P9, između obrazovanja i odgovora na P10 te između odgovora na P9 i P10. Najvažniji nalaz je povezanost između percipirane glavne prednosti SMR tehnologije i stava prema subvencioniranju njezina razvoja. Ispitanici koji jasno prepoznaju neku prednost SMR tehnologije, posebno sigurnosne značajke ili kraću izgradnju, češće podržavaju državne subvencije. Nasuprot tome, ispitanici koji nisu sigurni oko prednosti SMR tehnologije češće nisu sigurni ni oko opravdanosti subvencija.

Zaključno, rezultati upućuju na relativno povoljan stav ispitanika prema razvoju SMR tehnologije, ali i na važnost informiranja javnosti. Podrška javnom ulaganju u SMR tehnologiju povezana je s time prepoznaju li ispitanici konkretne prednosti te tehnologije. Stoga bi buduća komunikacija o SMR tehnologiji trebala jasno razdvojiti sigurnosne, ekonomske i energetske argumente te ih predstaviti na način koji omogućuje informirano oblikovanje javnog mnijenja.
